% ---------------------------------------------------------------------------
% A4SpectralRuntime.tex
% Altitude A4 — Spectral Runtime Physics
% Author: Borealis S. Hedling
% Version: v1.0
% ---------------------------------------------------------------------------

\section{Altitude A4: Spectral Runtime Physics}

Altitude A4 extends the Euclidean runtime manifold of A3 by introducing spectral
structure. Runtime operators acquire eigenfunctions, modal coordinates, and
spectral continuity constraints. A4 defines the mathematical transition from
adjacency-preserving Euclidean operators to spectrally-structured runtime
physics.

\subsection{Spectral Manifold Structure}

Let $M$ be the runtime manifold defined in A3. A4 introduces a spectral operator
$P : M \rightarrow M$ satisfying:
\[
P \text{ is linear, continuous, and adjacency-preserving}.
\]

A4 assumes the existence of a symmetric conjugate:
\[
P_{\mathrm{sym}} = D^{-1/2} K D^{-1/2},
\]
where $K$ is an adjacency-preserving kernel and $D$ is the degree matrix.

\subsection{Eigenfunctions and Spectral Modes}

A4 requires that $P_{\mathrm{sym}}$ admit an orthonormal eigenbasis:
\[
P_{\mathrm{sym}} \phi_k = \lambda_k \phi_k,
\]
with:
\[
1 = \lambda_0 \ge \lambda_1 \ge \lambda_2 \ge \cdots \ge 0.
\]

Eigenfunctions $\phi_k$ define the spectral modes of the runtime manifold.

\subsection{Modal Coordinates}

Modal coordinates are defined as:
\[
\psi_k = \lambda_k \phi_k.
\]

These coordinates embed runtime patches into a spectral space:
\[
\Psi : M \rightarrow \mathbb{R}^d, \quad p \mapsto (\psi_1(p), \ldots, \psi_d(p)).
\]

Modal coordinates must satisfy:
\begin{itemize}
  \item spectral continuity,
  \item adjacency preservation,
  \item bounded holonomy,
  \item reversible reconstruction.
\end{itemize}

\subsection{Spectral Continuity}

A4 enforces:
\[
\| \psi_k(p) - \psi_k(q) \| < \epsilon \quad \text{whenever} \quad A(p,q),
\]
ensuring that spectral modes do not fracture adjacency.

This condition prepares the manifold for A5 modal runtime.

\subsection{Spectral Stability}

Spectral stability requires:
\[
|\lambda_k - \lambda_k'| < \delta,
\]
under admissible perturbations of $K$ or $P$.

This prevents spectral tearing and ensures compatibility with A3 holonomy bounds.

\subsection{Spectral Operators}

An operator $T$ is admissible in A4 if:
\[
T \text{ preserves spectral continuity and modal coordinates}.
\]

Admissible operators include:
\begin{itemize}
  \item symmetric conjugate operators,
  \item Markov kernels,
  \item spectral filters,
  \item modal projections.
\end{itemize}

\subsection{DiffusionGeometry as a Spectral Example}

DiffusionGeometry provides spectral operators suitable for A4:
\begin{itemize}
  \item Markov operator $P$,
  \item symmetric conjugate $P_{\mathrm{sym}}$,
  \item eigenfunctions $\phi_k$,
  \item diffusion coordinates $\psi_k$.
\end{itemize}

These operators satisfy:
\[
K(p) = 0, \quad \|H_T\| < \epsilon, \quad \text{spectral continuity}.
\]

They remain strictly within altitudes A3–A4.

\subsection{Integration with Python Runtime Harness (A4.5)}

The Python runtime harness enforces:
\begin{itemize}
  \item spectral continuity constraints,
  \item adjacency-preserving spectral transitions,
  \item reversible modal coordinate updates,
  \item provenance monotonicity under spectral evolution.
\end{itemize}

A4SpectralRuntime.tex defines the spectral invariants the harness must satisfy.

% ---------------------------------------------------------------------------
% Provenance
% A4SpectralRuntime.tex
% Author: Borealis S. Hedling
% Altitude: A4
% Surface: Structural-AI/Runtime-Physics
% Version: v1.0
% Created: 18 October 2026, Dublin, Ireland
% Purpose: Define spectral runtime physics, eigenfunctions, modal coordinates,
%          spectral continuity, spectral stability, and admissible spectral
%          operators. Provides theoretical foundation for A5 modal runtime and
%          integrates with the Python A4.5 runtime harness.
% ---------------------------------------------------------------------------
